% !TeX root = ../../main.tex

PostgreSQL ist ein open-source Datenbankmanagementsystem (DBMS) das seit über 35 Jahren aktiv entwickelt wird.
Die aktuell veröffentlichte Version ist 18.4.
Version 19 wird gerade entwickelt und befindet sich in der Beta~\cite{PostgreSQL}.

Die Architektur von PostgreSQL sieht folgendermaßen aus.
Es gibt einen Serverprozess der die Datenbank verwaltet, Verbindungen akzeptiert und die Datenbankaktionen eines
Clients ausführt.
Dieses Programm heißt postgres.
Der Client kann dabei vielseitig sein, sei es ein textbasierter Client, eine grafische App oder ein Webserver.
Alle bauen eine TCP/IP-Verbindung zu dem Server auf, welcher mehrere Verbindungen gleichzeitig annehmen kann~\cite{PostgreSQLArchitecture}.

Für PostgreSQL Datenbanken sind Transaktionen das fundamentale Konzept.
Sie bündeln mehrere Schritte in eine einzige Operation.
Die Transaktionen basieren auf dem ACID Prinzip\footnote{Atomicity, eine Transaktion wird entweder ganz oder gar nicht ausgeführt.
	Consistency, nach jeder Transaktion bleiben alle Daten in einem gültigen und korrekten Zustand.
	Isolation, gleichzeitig ablaufende Transaktionen stören sich nicht gegenseitig.
	Durability, bestätigte Änderungen bleiben dauerhaft gespeichert~\cite{Haerder1983}.}, sodass eine Transaktion immer
gültig oder ungültig ist.
Eine Transaktion beginnt immer mit einem \texttt{BEGIN;} und endet mit einem \texttt{COMMIT;} oder \texttt{ROLLBACK;},
um eine Transaktion zu beenden~\cite{PostgreSQLTransactions}.

Um eine erfolgreiche Recovery bei Abstürzen zu gewährleisten, arbeitet PostgreSQL mit write ahead log (WAL)~\cite{PostgreSQLWAL}.
Klassisches WAL basiert auf \texttt{UNDO} Informationen, um unvollständige Transaktionen zurückzurollen.
Diese müssen vor der Propagation, \footnote{Propagation ist die Operation,
	die geschriebene Seiten zu einem Teil der materialized Database macht. Die materialized Database hingegen ist der
	persistente Datenbankzustand, welcher nach einem Absturz vorhanden ist~\cite{Haerder1983}.}
der entsprechenden Änderung in das Log geschrieben werden.
Dazu nutzt es \texttt{REDO} Informationen, um committete Änderungen zu wiederholen, die vor Bestätigung des
EOT\footnote{End of Transaction} in das aktive Log geschrieben werden müssen~\cite{Haerder1983}.
In PostgreSQL hingegen wird WAL anders verwendet.
Die Datendateien dürfen erst dann geschrieben werden, wenn die Änderungen protokolliert und auf einen permanenten
Speicher geschrieben wurden.
Dadurch muss nicht bei jedem Transaktionscommit auf die Festplatte geschrieben werden, denn im Falle eines Absturzes
kann die Datenbank mithilfe der Logs wiederhergestellt werden.
Alle Änderungen, die noch nicht propagiert wurden, können anhand der WAL-Einträge wiederholt werden, dies ist
der \texttt{REDO}~\cite{PostgreSQLWAL}.
